\documentclass[../../../../main.tex]{subfiles}
\begin{document}

\begin{flushright}
\footnotesize\textit{【自動文字起こし・要確認】}
\end{flushright}

\noindent
[\textbf{解}]
\[
\begin{cases}
2 \le a_1 \le a_2 \le \dots \le a_n \quad \dots \text{①} \\
S = n-2 - \displaystyle\sum_{i=1}^n \frac{1}{a_i} > 0
\end{cases}
\]

\begin{enumerate}
    \item[(i)] $n=3$ の時
    \[
    S = 1 - \left( \frac{1}{a_1} + \frac{1}{a_2} + \frac{1}{a_3} \right)
    \]
    だから、$S > 0$ をみたすうち、逆数和を max にするような $a$ の組をもとめれば良いが、$S$ は $2 \le x$ で単調減少なことから、$S$ ちいさくして考える。

    \begin{enumerate}
        \item[$1^\circ$] $a_1 = 2$ の時\\
        $\frac{1}{a_2} + \frac{1}{a_3} < \frac{1}{2}$ をみたすうち逆数和を max にする $a_2, a_3$ をかんがえる。$a_2, a_3 \ge 2$ から、$(a_2-2)(a_3-2) > 4$ だから、
        \[
        (a_2, a_3) = (3, 7), (4, 5) \text{ 又は } a_2 \ge 5
        \]
        なので、もとめる組は $(a_1, a_2, a_3) = (2, 3, 7)$。

        \item[$2^\circ$] $a_1 = 3$ の時\\
        $\frac{1}{a_2} + \frac{1}{a_3} < \frac{1}{3} \iff (a_2-3)(a_3-3) > 9$ をみたす逆数和を max にする $a_2, a_3$ をもとめる。$(a_2, a_3) = (4, 13), (5, 8), (6, 7)$ or $a_2 \ge 7$ で、明らかに $1^\circ$ の時より和が小さい。
    \end{enumerate}
    以上から、もとめるのは $(a_1, a_2, a_3) = (2, 3, 7)$。

    \item[(ii)] $n \ge 5$ の時、$a_i = 2\ (i=1, 2, \dots, n)$ とすることになり
    \[
    S = n-2 - \frac{n}{2} = \frac{1}{2}(n-4) > 0
    \]
    となって $S$ は最小値をとる。$n$ をうごかした時、これは明らかに $n=5$ で最小値をとる！以下 $n \le 4$ として考える。$n=1$ の時 $S$ は明らかに負だから不適。よって $n=4$ をしらべる。$\frac{1}{x}$ が $2 \le x$ で単調増加なことと $2 = a_1 = a_2 = a_3 = a_4$ の時 $S=0$ であることから、$a_1 = a_2 = a_3 = 2, a_4 = 3$ の時 $\min S = 2 - \left(\frac{1}{2} + \frac{1}{2} + \frac{1}{2} + \frac{1}{3}\right) = \frac{1}{6}\ (>0)$ をとる。
\end{enumerate}

\noindent
以上から、
\[
\begin{cases}
n \ge 5 \text{ の時 } n=5, a_i=2 \text{ で } \min S = \frac{1}{2} \\
n = 4 \text{ 〃 } a_1=a_2=a_3=2, a_4=3 \text{ で } \min S = \frac{1}{6} \\
n = 3 \text{ 〃 } (a_1, a_2, a_3)=(2,3,7) \text{ で } \min S = \frac{1}{42}
\end{cases}
\]
よって、$S$ を min にするのは
\[
n=3, \quad (a_1, a_2, a_3) = (2, 3, 7)
\]
の時。

\newpage

\end{document}
